% =============================================================================
\chapter{Várias placas, vários alunos}
\label{cap:placas}
% =============================================================================

O laboratório tem duas placas e uma turma. Este capítulo descreve o que
acontece entre um cliente e as placas quando várias pessoas usam o sistema ao
mesmo tempo: como o cliente acha as placas, como escolhe uma, o que a fila do
servidor faz com o tempo de resposta e o que foi medido. O texto completo, com
as mensagens de diagnóstico, está em \arq{docs/GERENCIAMENTO-PLACAS.md}.

\section{Descoberta}

O cliente procura placas em duas etapas.

\paragraph{A lista local.} O \arq{morphe-up.sh} grava toda placa que prepara
em \arq{.morphe-estado/placas}, um endereço por linha, na raiz do repositório;
a última preparada fica também em \arq{.morphe-estado/placa}. O cliente lê esses
arquivos primeiro. Como a instalação de turma é uma só (\arq{/opt/morphe}), o
que o \arq{morphe-up.sh} registra é o que todos os alunos daquele computador
enxergam.

\paragraph{A varredura.} Com a lista vazia, o cliente varre a rede por TCP,
enviando um \texttt{PING} a cada endereço das sub-redes prováveis --- a da placa
lembrada, a do próprio computador e as usadas historicamente no laboratório ---
com até 64 conexões em paralelo e prazos curtos. A varredura \textbf{para na
primeira placa que achar}. É o último recurso: o servidor não responde a
\emph{broadcast}, então uma placa numa sub-rede nunca vista continua invisível.

A lista é texto puro. Num computador que não é a estação, basta escrevê-la uma
vez para o cliente passar a conhecer as duas placas:

\begin{center}
\ttfamily printf '172.16.230.24\textbackslash n172.16.230.52\textbackslash n' > .morphe-estado/placas
\end{center}

\noindent Não se deve usar o \arq{morphe-up.sh} só para registrar uma placa:
mesmo sem reprogramar a FPGA, ele reinicia o servidor e interrompe quem estiver
no meio de uma operação.

\section{Escolha: a placa que responder mais rápido}

Com mais de uma placa na lista, o cliente envia um \texttt{PING} a todas ao
mesmo tempo e fica com a primeira que responder (Figura~\ref{fig:escolha-placa}).
Placas que respondem com \texttt{fpga\_preparada=0} são ignoradas.

\begin{figure}[htbp]
  \centering
  \figura{escolha-placa}
  \caption{Escolha da placa: o cliente sonda todas em paralelo e fica com a
    que responder primeiro. Uma placa ocupada só responde ao fim da operação
    em curso.}
  \label{fig:escolha-placa}
\end{figure}

O tempo de resposta do \texttt{PING} mede ocupação por causa da estrutura do
servidor: ele atende uma conexão de cada vez, então uma placa no meio de uma
convolução só responde ao \texttt{PING} depois de terminá-la. Placa livre
responde em milissegundos; placa ocupada, em centenas.

Duas placas livres empatam: a diferença entre elas é de microssegundos de rede,
e a vencedora muda de uma vez para a outra. Não é defeito. É também por isso
que dois alunos que cliquem no mesmo instante ainda podem cair na mesma placa;
a regra resolve o caso comum, de uma placa ocupada e outra livre.

\section{A fila e o tempo de resposta}

Como o servidor atende uma conexão por vez, dois alunos na mesma placa
\textbf{não recebem erro: esperam}. A segunda requisição fica na fila do sistema
operacional até a primeira terminar, e os tempos se somam. Com \(k\) clientes
disputando a mesma placa, cada um vê cerca de \(k\) vezes o tempo de placa
livre.

A fila do \texttt{listen(4)} guarda cinco conexões, e a sexta é a que está
sendo atendida. O que passa disso não é recusado: no Linux, o sistema descarta
o pedido de conexão em silêncio e o cliente o repete cerca de um segundo
depois. O sintoma do excesso é, portanto, uma cauda de latência, não um erro.

\section{O que foi medido}

As medições de 22/09/2026 usaram o \arq{Python/ferramentas/testa_concorrencia.py},
que dispara requisições simultâneas e \textbf{confere cada resposta}: uma
convolução com impulso tem de devolver a própria entrada, e a FFT de um impulso
tem de dar espectro plano. Não bastava cronometrar; era preciso saber se alguma
resposta chegava trocada.

\begin{table}[htbp]
  \centering
  \small
  \caption{Vários clientes simultâneos, operações de 1024 amostras, metade
    convoluções (\(\approx\)215\,ms) e metade FFTs (\(\approx\)21\,ms). Placa
    livre: mediana de 120\,ms (22/09/2026).}
  \label{tab:concorrencia}
  \begin{tabularx}{\textwidth}{@{}Xrrrr@{}}
    \toprule
    \textbf{Cenário} & \textbf{Respostas certas} & \textbf{Mediana}
      & \textbf{p95} & \textbf{Pior} \\
    \midrule
    4 clientes, uma placa  & 32/32 & 471\,ms (3,9\(\times\)) & --- & --- \\
    6 clientes, uma placa  & 24/24 & 604\,ms (5,1\(\times\)) & --- & 1,29\,s \\
    8 clientes, uma placa  & 64/64 & 517\,ms & 1,29\,s & 6,39\,s \\
    10 clientes, uma placa & 80/80 & 519\,ms & 4,56\,s & 7,07\,s \\
    \addlinespace
    3 clientes em cada uma de duas máquinas, duas placas & 90/90 & --- & --- & --- \\
    \bottomrule
  \end{tabularx}
\end{table}

Três conclusões:

\begin{itemize}
  \item \textbf{Nenhuma resposta foi perdida nem trocada}, em nenhum cenário. As
    90 requisições do teste com duas máquinas e duas placas foram conferidas uma
    a uma.
  \item \textbf{Uma placa atende cerca de 8 requisições por segundo} com essa
    mistura de operações, e esse número não muda de 2 para 10 clientes: é o
    tempo médio de uma operação (120\,ms) invertido. A partir de 8 clientes a
    placa está saturada, e cliente a mais só acrescenta espera. Duas placas dão
    cerca de 16 requisições por segundo para a turma inteira.
  \item O excesso aparece como cauda: de 8 para 10 clientes a mediana não se
    move e o p95 triplica.
\end{itemize}

Dois clientes gráficos, um em cada máquina, escolheram placas diferentes.

\paragraph{A captura do ADC ocupa a placa.} Uma captura contínua mantém a
conexão aberta enquanto dura, e durante esse tempo a placa não atende mais
ninguém: outro aluno na mesma placa espera, e a escolha automática não a
escolhe (o \texttt{PING} dela não responde). Numa aula com o ADC, é melhor
reservar uma placa para a aquisição.

\section{O que não existe}

Estas são as limitações conhecidas da infraestrutura. Nenhuma impede o uso em
turma com supervisão; todas impedem o uso sem supervisão com mais alunos do que
placas.

\begin{enumerate}
  \item \textbf{Alocação.} Nada impede duas pessoas de usarem a mesma placa
    achando que cada uma tem a sua. O resultado sai certo, porque cada
    requisição é completa, mas o tempo dobra e ninguém sabe por quê. Uma
    reserva com prazo, renovada pelo cliente, resolveria, e exige estado no
    servidor.
  \item \textbf{Concorrência no servidor.} Um processo por conexão, ou uma fila
    explícita que respondesse ``ocupada, você é o segundo'', diria ao cliente o
    que está acontecendo. Hoje ele não distingue ``placa ocupada'' de ``placa
    fora do ar'': nos dois casos a sonda simplesmente demora.
  \item \textbf{Descoberta sem varredura.} Um servidor que respondesse a
    \emph{broadcast} UDP seria achado em qualquer sub-rede.
  \item \textbf{Estado compartilhado entre computadores.} O
    \arq{.morphe-estado} é de cada instalação.
\end{enumerate}

\section{Sintoma, causa e o que fazer}

\begin{table}[htbp]
  \centering
  \small
  \caption{Sintomas ligados a várias placas e vários clientes.}
  \label{tab:sintomas-placas}
  \begin{tabularx}{\textwidth}{@{}XXX@{}}
    \toprule
    \textbf{Sintoma} & \textbf{Causa} & \textbf{O que fazer} \\
    \midrule
    A escolha automática alterna entre as placas & duas placas livres empatam
      & nada; é o esperado \\
    \addlinespace
    Tudo ficou duas vezes mais lento & outra pessoa está na mesma placa
      & é a fila; reconectar para ir à outra \\
    \addlinespace
    ``A placa não respondeu'', mas ela está no ar & placa ocupada (uma captura
      contínua, por exemplo) & esperar ou usar a outra \\
    \addlinespace
    Operação recusada com ``FPGA não preparada'' & a placa foi ligada e ninguém
      programou o bitstream & \arq{./morphe-up.sh} na estação \\
    \addlinespace
    Um computador só enxerga uma placa & a lista dele está vazia e a varredura
      para na primeira & escrever \arq{.morphe-estado/placas} \\
    \addlinespace
    ``o \emph{tether} de DE-SoC [1-x] pertence à conta \ldots'' & outra conta
      preparou aquela placa & usá-la como está ou pedir àquela conta um
      \opt{down} \\
    \bottomrule
  \end{tabularx}
\end{table}
